\documentclass[12pt,a4paper]{amsart}
\usepackage{amsmath,amssymb,amsthm}
\usepackage{mathtools}
\usepackage[utf8]{inputenc}
\usepackage[T1]{fontenc}
\usepackage{hyperref}
\usepackage{enumitem,booktabs,array}
\usepackage{geometry}
\geometry{margin=2.5cm}

\usepackage[ngerman]{babel}

\author{Michael Fuhrmann}
\address{Specialist Mathematics Project, Build 141}

\newtheorem{theorem}{Satz}[section]
\newtheorem{lemma}[theorem]{Lemma}
\newtheorem{korollar}[theorem]{Korollar}
\newtheorem{vermutung}[theorem]{Vermutung}
\theoremstyle{definition}
\newtheorem{definition}[theorem]{Definition}
\newtheorem{beispiel}[theorem]{Beispiel}
\theoremstyle{remark}
\newtheorem*{bemerkung}{Bemerkung}

\title{Zeta-Werte an ungeraden Stellen:\\
Irrationalitätsresultate, Apérys Satz und offene Probleme}

\begin{abstract}
Die Riemann-Zeta-Funktion $\zeta(s) = \sum_{n=1}^{\infty} n^{-s}$, ausgewertet
an positiven ganzzahligen Argumenten, beschäftigt Mathematiker seit Eulers Zeit.
Während die geraden Werte $\zeta(2n)$ vollständig verstanden sind --- sie sind
rationale Vielfache von Potenzen von $\pi$ --- bleibt die arithmetische Natur der
ungeraden Werte $\zeta(2n+1)$ weitgehend rätselhaft. Die einzige Ausnahme ist
$\zeta(3)$, dessen Irrationalität Roger Apéry im Jahr 1979 bewies und damit die
mathematische Gemeinschaft in Erstaunen versetzte. Wir präsentieren eine umfassende
Darstellung der Zeta-Werte an ungeraden ganzzahligen Stellen: Apérys Beweis der
Irrationalität von $\zeta(3)$, den Ball-Rivoal-Satz (2001), der unendlich viele
irrationale Werte unter $\{\zeta(2k+1)\}$ garantiert, und Zudilins Verschärfung
(2004), wonach mindestens einer von $\zeta(5), \zeta(7), \zeta(9), \zeta(11)$
irrational ist. Bernoulli-Zahlen, mehrfache Zeta-Werte, $p$-adische
Zeta-Funktionen und Verbindungen zu Modulformen werden ebenfalls behandelt.
Die Irrationalität von $\zeta(5)$ und die Transzendenz von $\zeta(3)$ bleiben
offene Vermutungen.
\end{abstract}

\maketitle

\tableofcontents

% ============================================================
\section{Einleitung}
% ============================================================

Die Riemann-Zeta-Funktion ist wohl die wichtigste Spezialfunktion der analytischen
Zahlentheorie. Für $\operatorname{Re}(s) > 1$ ist sie durch die absolut konvergente
Dirichlet-Reihe definiert:
\[
  \zeta(s) = \sum_{n=1}^{\infty} \frac{1}{n^s}.
\]
Sie setzt sich meromorph auf ganz $\mathbb{C}$ fort, mit einem einzigen einfachen
Pol bei $s = 1$ (Residuum~$1$). Eulers Produktformel
\[
  \zeta(s) = \prod_{p \text{ prim}} \frac{1}{1 - p^{-s}}, \quad
  \operatorname{Re}(s) > 1,
\]
verbindet die Zeta-Funktion mit der Verteilung der Primzahlen und gipfelt in der
Riemann-Hypothese, die besagt, dass alle nichttrivialen Nullstellen auf der
kritischen Geraden $\operatorname{Re}(s) = \tfrac{1}{2}$ liegen.

Über ihre Rolle in der Primzahltheorie hinaus wurden die Werte $\zeta(n)$ an
positiven ganzzahligen Argumenten intensiv untersucht. Euler bewies 1740, dass
\[
  \zeta(2n) = (-1)^{n+1} \frac{(2\pi)^{2n} B_{2n}}{2\,(2n)!}, \quad n \geq 1,
\]
wobei $B_{2n}$ die Bernoulli-Zahlen bezeichnet. Insbesondere gilt
$\zeta(2) = \pi^2/6$, $\zeta(4) = \pi^4/90$ usw. --- alles transzendente Zahlen,
die sich explizit durch $\pi$ ausdrücken lassen.

Bei den \emph{ungeraden} positiven ganzzahligen Argumenten ist die Lage
grundlegend verschieden. Keine analoge geschlossene Formel ist bekannt. Die Frage,
ob $\zeta(3), \zeta(5), \zeta(7), \ldots$ irrational, transzendent oder gar
algebraisch unabhängig über $\mathbb{Q}(\pi)$ sind, gehört zu den tiefsten offenen
Problemen der Zahlentheorie.

\subsection*{Historischer Überblick}

\begin{itemize}[leftmargin=2em]
  \item \textbf{Euler (1740)}: Bewies $\zeta(2n) \in \mathbb{Q} \cdot \pi^{2n}$.
        Der ungerade Fall blieb offen.
  \item \textbf{Apéry (1978/79)}: Bewies $\zeta(3) \notin \mathbb{Q}$ --- eine
        \emph{mathematische Sensation}, präsentiert bei den Journées Arithmétiques
        in Luminy (1978).
  \item \textbf{Beukers (1979)}: Gab eine elegante Umformulierung von Apérys
        Beweis mittels Legendre-Polynomen und Doppelintegralen.
  \item \textbf{Rivoal (2000) und Ball--Rivoal (2001)}: Bewiesen, dass unendlich
        viele ungerade Zeta-Werte irrational sind.
  \item \textbf{Zudilin (2004)}: Bewies, dass mindestens einer von
        $\zeta(5), \zeta(7), \zeta(9), \zeta(11)$ irrational ist.
\end{itemize}

Die Gliederung ist wie folgt: Abschnitt~2 behandelt den vollständig verstandenen
geraden Fall. Abschnitt~3 stellt Apérys Satz über $\zeta(3)$ im Detail vor.
Abschnitt~4 diskutiert die Konstante $\zeta(3)$ und ihre historische Bedeutung.
Abschnitt~5 behandelt die Sätze von Ball--Rivoal und Zudilin.
Abschnitt~6 führt mehrfache Zeta-Werte ein. Abschnitt~7 bespricht Bernoulli-Zahlen.
Abschnitt~8 behandelt $p$-adische Zeta-Funktionen. Abschnitt~9 verbindet die
Theorie mit Modulformen. Abschnitt~10 fasst offene Probleme zusammen.
Abschnitt~11 enthält die Literaturhinweise.

% ============================================================
\section{Gerade Werte: Eulers Formel}
% ============================================================

\subsection{Aussage und Beweisidee}

\begin{theorem}[Euler, 1740]
  Für jede positive ganze Zahl $n \geq 1$ gilt
  \[
    \zeta(2n) = (-1)^{n+1} \frac{(2\pi)^{2n} B_{2n}}{2\,(2n)!},
  \]
  wobei $B_{2n}$ die $2n$-te Bernoulli-Zahl bezeichnet.
\end{theorem}

\begin{proof}[Beweisidee]
  Der klassische Beweis benutzt die Produktentwicklung von $\sin(\pi z)$:
  \[
    \frac{\sin(\pi z)}{\pi z} = \prod_{n=1}^{\infty} \left(1 - \frac{z^2}{n^2}\right).
  \]
  Logarithmieren und Ableiten, oder Entwickeln beider Seiten als Potenzreihe in
  $z^2$ und Koeffizientenvergleich, liefert
  \[
    \sum_{n=1}^{\infty} \frac{1}{n^{2k}} = (-1)^{k+1} \frac{(2\pi)^{2k}
    B_{2k}}{2\,(2k)!}.
  \]
\end{proof}

\subsection{Tabelle der geraden Zeta-Werte}

\begin{table}[h]
  \centering
  \caption{Gerade Zeta-Werte $\zeta(2n)$ nach Eulers Formel.}
  \vspace{4pt}
  \begin{tabular}{@{}cclr@{}}
    \toprule
    $n$ & $\zeta(2n)$ (geschlossene Form) & Dezimalentwicklung & $B_{2n}$ \\
    \midrule
    1 & $\dfrac{\pi^2}{6}$           & $1{,}6449340668\ldots$ & $\tfrac{1}{6}$     \\[6pt]
    2 & $\dfrac{\pi^4}{90}$          & $1{,}0823232337\ldots$ & $-\tfrac{1}{30}$   \\[6pt]
    3 & $\dfrac{\pi^6}{945}$         & $1{,}0173430620\ldots$ & $\tfrac{1}{42}$    \\[6pt]
    4 & $\dfrac{\pi^8}{9450}$        & $1{,}0040773562\ldots$ & $-\tfrac{1}{30}$   \\[6pt]
    5 & $\dfrac{\pi^{10}}{93555}$    & $1{,}0009945751\ldots$ & $\tfrac{5}{66}$    \\[6pt]
    6 & $\dfrac{691\pi^{12}}{638512875}$ & $1{,}0002460866\ldots$ & $-\tfrac{691}{2730}$ \\
    \bottomrule
  \end{tabular}
  \label{tab:gerade-zeta}
\end{table}

\subsection{Folgerungen}

Da jedes $B_{2n} \in \mathbb{Q}$ gilt, ergibt Eulers Satz unmittelbar
$\zeta(2n) \in \mathbb{Q} \cdot \pi^{2n}$ für alle $n \geq 1$. Da $\pi$
transzendent ist (Lindemann 1882), ist jedes $\zeta(2n)$ transzendent.

Der \emph{ungerade} Fall offenbart die volle Tiefe des Problems: Es ist keine
Produktentwicklung bekannt, die $\zeta(2n+1)$ mit $\pi$ in Verbindung brächte.

% ============================================================
\section{Die Irrationalität von $\zeta(3)$: Apérys Satz}
% ============================================================

\subsection{Aussage}

\begin{theorem}[Apéry, 1979]\label{thm:apery}
  Die Zahl $\zeta(3) = \sum_{n=1}^{\infty} n^{-3}$ ist irrational.
\end{theorem}

Dieser Satz wurde bei den Journées Arithmétiques in Luminy im Juni 1978 angekündigt
und 1979 veröffentlicht. Das Resultat wurde zunächst skeptisch aufgenommen ---
mehrere Experten zweifelten am Beweis --- wurde aber innerhalb weniger Wochen von
Henri Cohen und anderen bestätigt. Bei der Ankündigung soll Apéry mit stehenden
Ovationen bedacht worden sein.

\subsection{Die Apéry-Folgen}

Die zentrale Idee von Apérys Beweis besteht darin, zwei Folgen positiver ganzer
Zahlen $\{a_n\}$ und $\{b_n\}$ zu konstruieren, die derselben dreigliedrigen
linearen Rekurrenz genügen und zu rationalen Approximationen von $\zeta(3)$
konvergieren --- so gute, dass $\zeta(3)$ nicht rational sein kann.

Die \emph{Apéry-Zahlen} sind definiert durch
\[
  b_n = \sum_{k=0}^{n} \binom{n}{k}^2 \binom{n+k}{k}^2, \quad n \geq 0.
\]
Die ersten Werte sind $b_0 = 1$, $b_1 = 5$, $b_2 = 73$, $b_3 = 1445$.

Die Zählerfolge $a_n$ ist definiert durch
\[
  a_n = \sum_{k=0}^{n} \binom{n}{k}^2 \binom{n+k}{k}^2
        \sum_{j=1}^{k} \frac{(-1)^{j-1}}{j^3 \binom{j+n}{j}\binom{n}{j}},
\]
so normiert, dass $a_n \in \mathbb{Z}$ nach Multiplikation mit einem geeigneten
kleinsten gemeinsamen Vielfachen.

\subsection{Die Schlüsselrekurrenz}

Beide Folgen genügen der \emph{Apéry-Rekurrenz}:
\begin{equation}
  (n+1)^3 u_{n+1} - (34n^3 + 51n^2 + 27n + 5)\, u_n + n^3 u_{n-1} = 0.
  \label{eq:apery-rekurrenz}
\end{equation}
Mit Anfangsbedingungen $(b_0, b_1) = (1, 5)$ entsteht die $b_n$-Folge; mit
$(a_0, a_1) = (0, 6)$ die Zählerfolge.

\subsection{Beweis der Irrationalität}

Die wesentlichen Abschätzungen sind:
\begin{enumerate}[label=(\roman*)]
  \item $b_n \zeta(3) - a_n \to 0$ für $n \to \infty$ (der Quotient konvergiert
        gegen $\zeta(3)$).
  \item Die Konvergenzrate ist
        \[
          |b_n \zeta(3) - a_n| \sim C \cdot (\sqrt{2} - 1)^{4n}, \quad C > 0,
        \]
        also superexponentiell schnell.
  \item Der Nenner $b_n$ wächst nur polynomiell: $b_n \sim A \cdot
        (1 + \sqrt{2})^{4n} / n^3$ für eine Konstante $A$.
\end{enumerate}

Zum Beweis der Irrationalität: Angenommen, $\zeta(3) = p/q$ für ganze Zahlen
$p, q > 0$. Dann gilt $|b_n \zeta(3) - a_n| = |b_n p/q - a_n| \geq 1/q$ für
große $n$ (als von null verschiedene ganze Zahl dividiert durch $q$), aber die
rechte Seite strebt schneller als jedes Polynom gegen null --- Widerspruch.

\subsection{Beukers' Integralformulierung}

Beukers (1979) gab eine elegante Umformulierung: Definiere
\begin{align*}
  I_n &= \int_0^1 \!\int_0^1 \frac{-\log(xy)}{1 - xy}\, P_n(x)\, P_n(y)\,
         \mathrm{d}x\, \mathrm{d}y, \\
  J_n &= \int_0^1 \!\int_0^1 \!\int_0^1
         \frac{x^n(1-x)^n\, y^n(1-y)^n\, z^n(1-z)^n}
              {(1 - (1-z)xy)^{n+1}}\,
         \mathrm{d}x\, \mathrm{d}y\, \mathrm{d}z,
\end{align*}
wobei $P_n$ das $n$-te Legendre-Polynom bezeichnet. Diese Integrale ergeben
$J_n = b_n \zeta(3) - a_n$ und liefern die Positivität und Ratenabschätzungen.

\subsection{Irrationalitätsmaß}

Die beste bekannte obere Schranke für das Irrationalitätsmaß von $\zeta(3)$
stammt von Rhin und Viola (2001):
\[
  \mu(\zeta(3)) \leq 5{,}5138905\ldots
\]
Das bedeutet: Für alle $\varepsilon > 0$ hat $|\zeta(3) - p/q| < q^{-5{,}5138905-\varepsilon}$
nur endlich viele Lösungen.

% ============================================================
\section{Apérys Konstante $\zeta(3)$}
% ============================================================

\subsection{Zahlenwert und Kettenbruch}

Die Konstante $\zeta(3)$ trägt heute den Namen \emph{Apérysche Konstante}:
\[
  \zeta(3) = 1{,}2020569031595942854523691932819321312725\ldots
\]
Ihre Kettenbruchentwicklung beginnt mit
\[
  \zeta(3) = [1;\, 4,\, 1,\, 18,\, 1,\, 1,\, 1,\, 4,\, 1,\, 9,\, 9,\, 2,\,
  1,\, 1,\, 1,\, 2, \ldots],
\]
ohne erkennbares Muster. Das stimmt mit der vermuteten (aber unbewiesenen)
Transzendenz überein.

\subsection{Alternative Reihendarstellungen}

Mehrere schneller konvergierende Reihen sind bekannt:

\begin{itemize}[leftmargin=2em]
  \item \textbf{Eulers Formel:}
    \[
      \zeta(3) = \frac{5}{2} \sum_{n=1}^{\infty} \frac{(-1)^{n+1}}{n^3 \binom{2n}{n}}.
    \]
  \item \textbf{Alternierende Dirichlet-Reihe:}
    \[
      \zeta(3) = \frac{4}{3} \sum_{n=1}^{\infty} \frac{(-1)^{n+1}}{n^3}.
    \]
  \item \textbf{Borwein--Crandall--Fee-Formel:}
    \[
      \zeta(3) = \frac{\pi^2}{7} \left( 1 - 4 \sum_{k=1}^{\infty}
      \frac{\zeta(2k)}{(2k+1)(2k+2)\, 4^k} \right).
    \]
\end{itemize}

\subsection{Apérys Wunder}

Der Begriff \emph{Apérys Wunder} (engl. \emph{Apéry's miracle}) bezieht sich auf
die erstaunliche Tatsache, dass die Koeffizienten $b_n$ in Apérys Beweis ganze
Zahlen sind --- was aus den kombinatorischen Summen \emph{a priori} keineswegs
offensichtlich ist. Die Ganzzahligkeit folgt aus Teilbarkeitseigenschaften der
Binomialkoeffizienten, hat aber keine einfache konzeptionelle Erklärung.
Alfred van der Poortens Bericht ``A proof that Euler missed'' (1979) machte
Apérys Argumentation einem breiten Publikum zugänglich.

\subsection{Offener Status}

\begin{vermutung}
  $\zeta(3)$ ist transzendent (d.\,h. nicht algebraisch über $\mathbb{Q}$).
\end{vermutung}

Dies bleibt vollständig offen (Stand 2026). Den Beweis der Irrationalität zu
erbringen war bereits schwierig; Transzendenz erfordert völlig neue Methoden.

% ============================================================
\section{Irrationalität unter $\zeta(5), \zeta(7), \ldots$}
% ============================================================

\subsection{Der Satz von Ball und Rivoal}

\begin{theorem}[Ball--Rivoal, 2001]\label{thm:ball-rivoal}
  Die $\mathbb{Q}$-Vektorraumdimension von
  \[
    \operatorname{span}_{\mathbb{Q}}\{1,\, \zeta(3),\, \zeta(5),\, \ldots,\,
    \zeta(2n+1)\}
  \]
  ist mindestens
  \[
    \frac{(1 + o(1))\,\ln n}{1 + \ln 2}
    \quad \text{für } n \to \infty.
  \]
  Insbesondere sind unendlich viele der Werte $\zeta(2k+1)$ irrational.
\end{theorem}

\begin{proof}[Beweisidee]
  Der Beweis konstruiert für jedes große $r$ eine Linearform
  \[
    L_r = c_0^{(r)} + c_3^{(r)}\,\zeta(3) + c_5^{(r)}\,\zeta(5) + \cdots
          + c_{2r-1}^{(r)}\,\zeta(2r-1)
  \]
  mit ganzzahligen Koeffizienten $c_j^{(r)}$, wobei $L_r \neq 0$, aber
  $|L_r| \to 0$. Die Koeffizienten entstammen einem hypergeometrischen Integral,
  und ihre Kleinheit relativ zu den Nennern erzwingt, dass die Menge
  $\{1, \zeta(3), \ldots, \zeta(2r-1)\}$ mindestens $\sim \ln(r)/(1 + \ln 2)$
  linear unabhängige Elemente über $\mathbb{Q}$ enthält.
\end{proof}

\subsection{Die Wachstumsrate}

Genauer gilt
\[
  \dim_{\mathbb{Q}} \operatorname{span}_{\mathbb{Q}}\{1, \zeta(3), \zeta(5),
  \ldots, \zeta(2n+1)\} \geq \frac{\ln n}{1 + \ln 2} \approx \frac{\ln n}{1{,}6931}
\]
für große $n$. Die untere Schranke wächst logarithmisch.

\subsection{Zudilins Satz}

\begin{theorem}[Zudilin, 2004]\label{thm:zudilin}
  Mindestens einer der vier Werte
  \[
    \zeta(5),\quad \zeta(7),\quad \zeta(9),\quad \zeta(11)
  \]
  ist irrational.
\end{theorem}

\begin{proof}[Beweisidee]
  Zudilin konstruierte eine Linearform
  \[
    L = a_0 + a_5 \zeta(5) + a_7 \zeta(7) + a_9 \zeta(9) + a_{11} \zeta(11),
    \quad a_i \in \mathbb{Z},
  \]
  mit $L \neq 0$ aber $|L| < 1$. Wären alle vier Werte rational mit einem
  gemeinsamen Nenner $D$, so wäre $D \cdot L$ eine von null verschiedene ganze
  Zahl in $(-1, 1)$ --- Widerspruch.

  Die Konstruktion verwendet hypergeometrische Funktionen ${}_{7}F_6$ und das
  Siegel-Lemma, um geeignete ganzzahlige Koeffizienten mit kontrollierter Größe
  zu finden.
\end{proof}

\begin{bemerkung}
  Zudilins Satz identifiziert \emph{nicht}, welcher der vier Werte irrational ist.
  Die Irrationalität von $\zeta(5)$ selbst bleibt eine offene Vermutung.
\end{bemerkung}

\subsection{Tabelle der ungeraden Zeta-Werte}

\begin{table}[h]
  \centering
  \caption{Ungerade Zeta-Werte und ihr arithmetischer Status.}
  \vspace{4pt}
  \begin{tabular}{@{}clll@{}}
    \toprule
    $n$ & $\zeta(n)$ (Näherung) & Irrationalität & Quelle \\
    \midrule
    3  & $1{,}20205690315959\ldots$ & \textbf{Bewiesen irrational} & Apéry (1979) \\
    5  & $1{,}03692775514337\ldots$ & Offene Vermutung            & --- \\
    7  & $1{,}00834927738192\ldots$ & Offene Vermutung            & --- \\
    9  & $1{,}00200839292608\ldots$ & Offene Vermutung            & --- \\
    11 & $1{,}00049418860409\ldots$ & Offene Vermutung            & --- \\
    13 & $1{,}00012271334757\ldots$ & Offene Vermutung            & --- \\
    15 & $1{,}00003058823630\ldots$ & Offene Vermutung            & --- \\
    \midrule
    \multicolumn{4}{l}{Mindestens einer von $\zeta(5), \zeta(7), \zeta(9), \zeta(11)$
    irrational: \textbf{Bewiesen} (Zudilin 2004)} \\
    \multicolumn{4}{l}{Unendlich viele $\zeta(2k+1)$ irrational: \textbf{Bewiesen}
    (Ball--Rivoal 2001)} \\
    \bottomrule
  \end{tabular}
  \label{tab:ungerade-zeta}
\end{table}

% ============================================================
\section{Mehrfache Zeta-Werte}
% ============================================================

\subsection{Definition}

Mehrfache Zeta-Werte (engl. \emph{multiple zeta values}, MZV) verallgemeinern
die Riemann-Zeta-Funktion auf Multiindex-Argumente. Für ganze Zahlen $s_1 \geq 2$
und $s_2, \ldots, s_k \geq 1$ setzt man
\[
  \zeta(s_1, s_2, \ldots, s_k) = \sum_{n_1 > n_2 > \cdots > n_k \geq 1}
  \frac{1}{n_1^{s_1} n_2^{s_2} \cdots n_k^{s_k}}.
\]
Die ganze Zahl $k$ heißt \emph{Tiefe}, die Summe $|s| = s_1 + \cdots + s_k$
heißt \emph{Gewicht}.

\subsection{Beispiele und erste Werte}

\begin{align*}
  \zeta(2) &= \frac{\pi^2}{6}, \\
  \zeta(3) &= 1{,}20205690\ldots \quad \text{(Apérysche Konstante)}, \\
  \zeta(2, 1) &= \zeta(3) \quad \text{(Euler-Resultat)}, \\
  \zeta(1, 2) &= \zeta(3) \quad \text{(Dualität)}, \\
  \zeta(2, 3) &= \tfrac{9}{2}\,\zeta(5) - 2\,\zeta(2)\zeta(3).
\end{align*}

\subsection{Die Algebra der MZV}

MZV genügen zwei Familien algebraischer Relationen:

\textbf{Shuffle-Produkt.} Partielle Integration ergibt
\[
  \zeta(\mathbf{s}) \cdot \zeta(\mathbf{t}) = \sum_{\sigma \in
  \mathrm{Sh}(\mathbf{s}, \mathbf{t})} \zeta(\sigma),
\]
wobei die Summe über alle Shuffle-Permutationen läuft.

\textbf{Stuffle-Produkt (Quasi-Shuffle).} Die Multiplikation von Dirichlet-Reihen
liefert eine andere Identität:
\[
  \zeta(s_1) \cdot \zeta(s_2) = \zeta(s_1, s_2) + \zeta(s_2, s_1) +
  \zeta(s_1 + s_2).
\]

Diese beiden Produktstrukturen erzwingen starke lineare Relationen unter MZV.
Die \emph{Doppel-Shuffle-Relationen} erzeugen alle bekannten algebraischen
Relationen, und es wird vermutet (aber nicht bewiesen), dass sie
\emph{alle} Relationen erzeugen.

\subsection{Gewicht und Dimension}

MZV vom Gewicht $w$ spannen vermutlich einen $\mathbb{Q}$-Vektorraum der Dimension
$d_w$ auf, wobei $d_w = d_{w-2} + d_{w-3}$ (Zagiers Vermutung). Die ersten Werte
sind:
\[
  d_1 = 0,\quad d_2 = 1,\quad d_3 = 1,\quad d_4 = 1,\quad d_5 = 2,
  \quad d_6 = 2, \quad d_7 = 3.
\]

% ============================================================
\section{Bernoulli-Zahlen und Euler--Maclaurin}
% ============================================================

\subsection{Definition und Rekurrenz}

Die Bernoulli-Zahlen $B_n$ sind durch die erzeugende Funktion definiert:
\[
  \frac{t}{e^t - 1} = \sum_{n=0}^{\infty} B_n \frac{t^n}{n!}, \quad |t| < 2\pi.
\]
Sie genügen der Rekurrenz
\[
  \sum_{k=0}^{n} \binom{n+1}{k} B_k = 0, \quad n \geq 1,
\]
mit $B_0 = 1$. Es gilt $B_1 = -1/2$, $B_{2k+1} = 0$ für $k \geq 1$, und

\begin{table}[h]
  \centering
  \caption{Bernoulli-Zahlen $B_{2k}$.}
  \vspace{4pt}
  \begin{tabular}{@{}crr@{}}
    \toprule
    $k$ & $B_{2k}$ & Dezimalwert \\
    \midrule
    1 & $1/6$               & $0{,}1\overline{6}$ \\
    2 & $-1/30$             & $-0{,}0\overline{3}$ \\
    3 & $1/42$              & $0{,}0238\ldots$ \\
    4 & $-1/30$             & $-0{,}0\overline{3}$ \\
    5 & $5/66$              & $0{,}07575\ldots$ \\
    6 & $-691/2730$         & $-0{,}25311\ldots$ \\
    7 & $7/6$               & $1{,}1\overline{6}$ \\
    \bottomrule
  \end{tabular}
\end{table}

\subsection{Euler--Maclaurinsche Summenformel}

\begin{theorem}[Euler--Maclaurin]
  Für eine glatte Funktion $f$ auf $[1, N]$ gilt
  \[
    \sum_{n=1}^{N} f(n) = \int_1^N f(x)\,\mathrm{d}x
    + \frac{f(1) + f(N)}{2}
    + \sum_{k=1}^{p} \frac{B_{2k}}{(2k)!}\bigl(f^{(2k-1)}(N) - f^{(2k-1)}(1)\bigr)
    + R_p,
  \]
  wobei $R_p$ ein Rest mit $f^{(2p+1)}$ ist.
\end{theorem}

Angewendet auf $f(x) = x^{-s}$ liefert die Euler--Maclaurin-Formel die analytische
Fortsetzung von $\zeta(s)$ und für negative ganze Zahlen die Spezialwerte
$\zeta(-n) = -B_{n+1}/(n+1)$.

\subsection{Kummers Kongruenzen}

Für eine Primzahl $p$ und $p \nmid k$ besagen Kummers Kongruenzen:
\[
  \frac{B_k}{k} \equiv \frac{B_{k'}}{k'} \pmod{p}
\]
wann immer $k \equiv k' \pmod{(p-1)}$ und $p \nmid k$.
Dies ist der Eckpfeiler der $p$-adischen Interpolation von Zeta-Werten.

% ============================================================
\section{$p$-adische Zeta-Funktionen}
% ============================================================

\subsection{Motivation}

Die rationalen Werte $\zeta(1-n) = -B_n/n$ für positive ganze Zahlen $n$ genügen
Kummers Kongruenzen, was nahelegt, dass $n \mapsto \zeta(1-n)$ stetig auf
$p$-adische ganze Zahlen ausgedehnt werden kann.

\subsection{Kubota--Leopoldt $p$-adische $L$-Funktionen}

\begin{theorem}[Kubota--Leopoldt, 1964]
  Für jede Primzahl $p$ und jeden Dirichlet-Charakter $\chi$ mit $\chi(-1) = (-1)^k$
  existiert eine eindeutige stetige $p$-adische Funktion
  \[
    L_p(s, \chi): \mathbb{Z}_p \to \mathbb{C}_p,
  \]
  so dass für alle positiven ganzen Zahlen $k$ mit $k \equiv 0 \pmod{p-1}$
  (bzw. $k$ gerade im Fall $\chi = \mathbf{1}$) gilt:
  \[
    L_p(1-k, \chi) = -\frac{B_{k,\chi}}{k} \cdot \left(1 - \chi(p) p^{k-1}\right),
  \]
  wobei $B_{k,\chi}$ die verallgemeinerte Bernoulli-Zahl ist.
\end{theorem}

\subsection{$p$-adische Interpolation und Iwasawa-Theorie}

Die Kubota--Leopoldt-Funktion ist das zentrale Objekt der Iwasawa-Theorie. In der
zyklotomischen $\mathbb{Z}_p$-Erweiterung $\mathbb{Q}_\infty$ von $\mathbb{Q}$
liefert der projektive Limes der Klassengruppen einen endlich erzeugten
$\Lambda$-Modul (wobei $\Lambda = \mathbb{Z}_p[[T]]$ die Iwasawa-Algebra ist),
dessen charakteristisches Polynom eng mit $L_p(s, \chi)$ zusammenhängt.

Die \emph{Hauptvermutung der Iwasawa-Theorie} (von Mazur--Wiles 1984 für $\mathbb{Q}$
mit Modulkurven bewiesen) besagt, dass das charakteristische Ideal dieses Moduls
gleich dem von der $p$-adischen $L$-Funktion erzeugten Ideal ist.

\subsection{Zusammenhang mit ungeraden Zeta-Werten}

Die $p$-adische Zeta-Funktion interpoliert die Werte $\zeta(1-k)$ für gerades
positives $k$, nicht direkt die ungeraden positiven Werte. Das Fehlen einer
$p$-adischen Interpolation für ungerade positive ganze Zahlen ist selbst eine
Manifestation des arithmetischen Rätsels um $\zeta(2k+1)$.

% ============================================================
\section{Verbindung zu Modulformen}
% ============================================================

\subsection{Eisenstein-Reihen und $\zeta(2k)$}

Die Eisenstein-Reihe vom Gewicht $2k$ für $\mathrm{SL}_2(\mathbb{Z})$ ist
\[
  G_{2k}(\tau) = \sum_{\substack{(m,n) \in \mathbb{Z}^2 \\ (m,n) \neq (0,0)}}
  \frac{1}{(m\tau + n)^{2k}}, \quad k \geq 2, \quad \tau \in \mathbb{H}.
\]
Ihre Fourier-Entwicklung lautet
\[
  G_{2k}(\tau) = 2\zeta(2k) + \frac{2(2\pi i)^{2k}}{(2k-1)!}
  \sum_{n=1}^{\infty} \sigma_{2k-1}(n)\, q^n, \quad q = e^{2\pi i \tau},
\]
wobei $\sigma_r(n) = \sum_{d|n} d^r$. Dies drückt $\zeta(2k)$ direkt als
konstanten Term einer Eisenstein-Reihe aus und erklärt das Auftreten von $\pi^{2k}$.

\subsection{Warum ungerade Werte anders sind}

Für $\mathrm{SL}_2(\mathbb{Z})$ gibt es keine Eisenstein-Reihen ungeraden
Gewichts (da $\mathrm{SL}_2(\mathbb{Z})$ das Element $-I$ enthält, das Modulformen
ungeraden Gewichts zum Verschwinden zwingt). Dies ist ein tiefer struktureller
Grund, warum $\zeta(2k+1)$ keinen einfachen Ausdruck durch $\pi$ besitzt.

\subsection{Die Ramanujan--Berndt-Formel für $\zeta(2k+1)$}

Ramanujan (und später Berndt) bewiesen die Formel
\[
  \zeta(2k+1) = \frac{2^{2k} \pi^{2k+1}}{(2k+1)!}
  \left[
    - B_{2k+1} + 2 \sum_{n=1}^{\infty} \frac{\sigma_{2k+1}^*(n)}{e^{2\pi n} - 1}
  \right],
\]
wobei $\sigma_r^*$ Teilersummen enthält und $B_{2k+1}$ verallgemeinerte
Bernoulli-Zahlen sind. Diese Formel reduziert $\zeta(2k+1)$ jedoch
\emph{nicht} auf ein rationales Vielfaches von $\pi^{2k+1}$.

\subsection{Perioden-Integrale}

In der Sprache der Motive sind gerade Zeta-Werte $\zeta(2n)$ \emph{Perioden
gemischter Tate-Motive}, die über $\mathbb{Q}$ splitten, während ungerade
Zeta-Werte mit der \emph{motivischen Kohomologie} von $\mathrm{Spec}(\mathbb{Z})$
zusammenhängen. Dieser von Beilinson, Deligne und Goncharov entwickelte Rahmen
bietet die tiefste Perspektive auf die Arithmetik der Zeta-Werte.

% ============================================================
\section{Offene Probleme}
% ============================================================

\begin{vermutung}[Transzendenz von $\zeta(3)$]
  Die Konstante $\zeta(3) = 1{,}2020569\ldots$ ist transzendent.
\end{vermutung}

\begin{vermutung}[Irrationalität von $\zeta(5)$]
  Der Wert $\zeta(5) = 1{,}0369277\ldots$ ist irrational.
\end{vermutung}

\begin{vermutung}[Irrationalität aller ungeraden Zeta-Werte]
  Für jedes $k \geq 1$ ist der Wert $\zeta(2k+1)$ irrational.
\end{vermutung}

\begin{vermutung}[$\zeta(2n+1)/\pi^{2n+1}$ irrational]
  Für jedes $k \geq 1$ ist der Quotient $\zeta(2k+1)/\pi^{2k+1}$ irrational
  (und in der Tat nicht algebraisch).
\end{vermutung}

\begin{vermutung}[Algebraische Unabhängigkeit]
  Die Zahlen $\pi, \zeta(3), \zeta(5), \zeta(7), \ldots$ sind algebraisch
  unabhängig über $\mathbb{Q}$.
\end{vermutung}

\subsection*{Statusübersicht}

\begin{table}[h]
  \centering
  \caption{Offene Probleme zu ungeraden Zeta-Werten (Stand 2026).}
  \vspace{4pt}
  \begin{tabular}{@{}lll@{}}
    \toprule
    Behauptung & Status & Quelle \\
    \midrule
    $\zeta(3)$ irrational & \textbf{BEWIESEN} & Apéry (1979) \\
    $\zeta(2k+1)$ irrational ($\infty$ viele) & \textbf{BEWIESEN} & Ball--Rivoal (2001) \\
    Mind. einer von $\zeta(5),\ldots,\zeta(11)$ irrational & \textbf{BEWIESEN} & Zudilin (2004) \\
    $\zeta(3)$ transzendent & Offene Vermutung & --- \\
    $\zeta(5)$ irrational & Offene Vermutung & --- \\
    Alle $\zeta(2k+1)$ irrational & Offene Vermutung & --- \\
    $\zeta(2k+1)/\pi^{2k+1}$ irrational & Offene Vermutung & --- \\
    Alg. Unabhängigkeit von $\pi, \zeta(3), \zeta(5), \ldots$ & Offene Vermutung & --- \\
    \bottomrule
  \end{tabular}
\end{table}

\subsection*{Aussichten und Ansätze}

Wesentlicher Fortschritt würde voraussichtlich einen oder mehrere der folgenden
Schritte erfordern:
\begin{enumerate}[leftmargin=2em]
  \item Ein besseres Verständnis der \emph{motivischen Galois-Gruppe} und ihrer
        Wirkung auf Perioden.
  \item Neue hypergeometrische Identitäten, die kleinere Linearformen in $\zeta(5)$
        liefern.
  \item Eine $p$-adische Interpolation ungerader Zeta-Werte (die neue Strukturen
        über Kubota--Leopoldt hinaus erfordern würde).
  \item Effektive Versionen des Lindemann--Weierstrass-Satzes für Polylogarithmen.
\end{enumerate}

% ============================================================
\section{Literatur}
% ============================================================

\begin{thebibliography}{99}

\bibitem{Apery1979}
R.~Apéry,
\textit{Irrationalité de $\zeta(2)$ et $\zeta(3)$},
Astérisque \textbf{61} (1979), 11--13.

\bibitem{BallRivoal2001}
K.~Ball und T.~Rivoal,
\textit{Irrationalité d'une infinité de valeurs de la fonction zeta aux entiers impairs},
Inventiones Mathematicae \textbf{146} (2001), 193--207.

\bibitem{Beukers1979}
F.~Beukers,
\textit{A note on the irrationality of $\zeta(2)$ and $\zeta(3)$},
Bulletin of the London Mathematical Society \textbf{11} (1979), 268--272.

\bibitem{Euler1740}
L.~Euler,
\textit{De summis serierum reciprocarum},
Commentarii Academiae Scientiarum Petropolitanae \textbf{7} (1740), 123--134.

\bibitem{RhinViola2001}
G.~Rhin und C.~Viola,
\textit{The group structure for $\zeta(3)$},
Acta Arithmetica \textbf{97} (2001), 269--293.

\bibitem{Rivoal2000}
T.~Rivoal,
\textit{La fonction zêta de Riemann prend une infinité de valeurs irrationnelles
aux entiers impairs},
Comptes Rendus de l'Académie des Sciences (Paris), Série~I \textbf{331} (2000),
267--270.

\bibitem{Zudilin2004}
W.~Zudilin,
\textit{Arithmetic of linear forms involving odd zeta values},
Journal de Théorie des Nombres de Bordeaux \textbf{16} (2004), 251--291.

\bibitem{vanderPoorten1979}
A.~van der Poorten,
\textit{A proof that Euler missed \ldots\ Apéry's proof of the irrationality of
$\zeta(3)$},
The Mathematical Intelligencer \textbf{1} (1979), 195--203.

\bibitem{Zagier1994}
D.~Zagier,
\textit{Values of zeta functions and their applications},
in: \textit{First European Congress of Mathematics}, Bd.~II,
Birkhäuser, Basel, 1994, S.~497--512.

\bibitem{KubotaLeopoldt1964}
T.~Kubota und H.~Leopoldt,
\textit{Eine $p$-adische Theorie der Zetawerte I},
Journal für die reine und angewandte Mathematik \textbf{214/215} (1964), 328--339.

\end{thebibliography}

\end{document}
